\chapter{Introduction}

Machine learning has transformed how computers solve complex tasks, from recognizing objects in images to generating human-like text. Most machine learning systems learn by studying large datasets of examples. Reinforcement learning (RL) takes a fundamentally different approach: instead of learning from labeled data, an RL agent learns by interacting with an environment through trial and error. At each step, the agent observes the current state of the environment, takes an action, and receives a reward signal indicating how good or bad that action was. Over many interactions, the agent learns a \emph{policy}, a strategy that maps observations to actions, that maximizes the cumulative reward it receives over time.

Training an RL agent involves two alternating phases. In the \emph{rollout phase}, the agent interacts with the environment using its current policy, collecting a batch of experiences consisting of observations, actions, and rewards. In the \emph{training phase}, the agent uses this collected experience to update its policy through gradient-based optimization. The agent's policy is represented by a neural network, and training adjusts the network's parameters to increase the probability of actions that led to higher rewards. This cycle of collecting experience and updating the policy repeats until the agent achieves satisfactory performance.

Modern RL agents often use deep neural networks with millions of parameters, particularly when the input is high-dimensional. For example, an agent learning to drive a car from camera images needs a convolutional neural network (CNN) to process visual observations before a policy network can select actions. As these networks grow in size and complexity, the computational and memory requirements for training can exceed what a single device can provide.

\section{Distributed Training}

Distributed training addresses the resource limitations of a single device by spreading computation across multiple machines. The simplest approach is \emph{centralized training}: an edge device such as a robot or sensor collects observations from the environment and sends them to a powerful cloud server, which runs the entire neural network and returns actions or updated model parameters. This works well when the model fits in the cloud server's memory and network bandwidth is sufficient for transmitting raw observations.

However, when a model is too large for a single device's memory, or when we want to keep parts of the computation close to the data source, we need a different strategy. \emph{Pipeline parallelism} partitions the neural network into sequential stages, with each stage running on a separate device. For example, a CNN feature extractor can run on one machine while the policy and value networks run on another. During the forward pass, the first machine processes raw observations into compact feature vectors and sends them to the second machine, which computes actions. During the backward pass for training, gradients flow in the reverse direction.

Pipeline parallelism introduces a communication challenge that centralized training does not face. In centralized training, the entire neural network resides on one machine, so all gradient computation and backpropagation happen locally in memory. The only network communication is the transfer of observations from the edge device to the cloud. In pipeline parallelism, the network split forces activations and gradients to be exchanged between machines for every minibatch during every training pass over the data. Since RL algorithms such as Proximal Policy Optimization (PPO)~\cite{ppo2017} reuse collected experience across multiple epochs for sample efficiency, this exchange is repeated many times per iteration, amplifying the communication overhead significantly.

\section{Communication Bottleneck}

Consider a concrete example. An RL agent learning from $96 \times 96$ RGB images with a batch of 4,096 steps per iteration and 10 update epochs of 32 minibatches each must exchange activations and gradients $10 \times 32 = 320$ times per iteration. Even after compressing observations through a CNN, the repeated gradient transmissions can exceed the total data volume of simply sending raw observations to a cloud server, as we quantify in Chapter~3. In bandwidth-constrained environments such as edge deployments or IoT networks, this communication overhead becomes the dominant bottleneck, slowing down training or making pipeline parallelism impractical altogether.

\section{Our Approach}

This project addresses the communication bottleneck through selective gradient transmission. Rather than sending all gradients from one machine to the other after every minibatch, we accumulate gradients locally and transmit only the most significant values based on magnitude percentile thresholds. Gradients that fall below the threshold are retained in a local buffer and continue accumulating until they become significant enough to send.

A key challenge with this accumulation approach is gradient staleness. Gradients that remain in the buffer across multiple iterations may no longer reflect the current optimization landscape by the time they are finally transmitted. To address this, we introduce an exponential decay factor that progressively discounts older unsent gradients, ensuring the buffer reflects recent training dynamics rather than outdated signals.

\section{Thesis Organization}

The remainder of this thesis is organized as follows. Chapter~2 reviews related work in distributed reinforcement learning, communication-efficient training, and pipeline parallelism. Chapter~3 defines the problem formally and quantifies the communication overhead of pipeline parallelism compared to centralized training. Chapter~4 describes our gradient accumulation method with percentile-based sparsification and buffer decay. Chapter~5 presents experimental results on the CarRacing-v3 environment, demonstrating a 35\% reduction in network transfer while maintaining comparable policy performance, and analyzes the effect of buffer decay on gradient accumulation.
